%======================================================================
%  notation.tex -- one place for the notation of the whole book
%
%  scalars ............ italic                 u, \theta, \lambda
%  vectors ............ bold italic            \vect{u}, \vu, \vn
%  2nd-order tensors .. bold italic            \tens{\sigma}, \sig, \eps
%  4th-order tensors .. blackboard bold        \tensf{C}, \bbC, \bbS, \bbI
%  matrices (discrete)  bold upright           \mat{K}, \mat{B}
%  column of unknowns . bold upright           \mat{u}  (= [u] in the slides)
%======================================================================
\newcommand{\vect}[1]{\bm{#1}}
\newcommand{\tens}[1]{\bm{#1}}
\newcommand{\tensf}[1]{\mathbb{#1}}
\newcommand{\mat}[1]{\mathbf{#1}}

% sets and spaces
\newcommand{\R}{\mathbb{R}}
\newcommand{\Vsp}{\mathcal{V}}          % functional space of the variational problem
\newcommand{\Kad}{\mathcal{K}}          % kinematically admissible fields
\newcommand{\Sad}{\mathcal{S}}          % statically admissible fields

% vectors
\newcommand{\vu}{\vect{u}}
\newcommand{\vv}{\vect{v}}
\newcommand{\vw}{\vect{w}}
\newcommand{\vx}{\vect{x}}
\newcommand{\vy}{\vect{y}}
\newcommand{\vn}{\vect{n}}
\newcommand{\vt}{\vect{t}}
\newcommand{\vf}{\vect{f}}
\newcommand{\vb}{\vect{b}}
\newcommand{\ve}{\vect{e}}
\newcommand{\vd}{\vect{d}}
\newcommand{\vs}{\vect{s}}
\newcommand{\vq}{\vect{q}}
\newcommand{\vp}{\vect{p}}
\newcommand{\vF}{\vect{F}}
\newcommand{\vE}{\vect{E}}
\newcommand{\vm}{\vect{m}}
\newcommand{\vxi}{\vect{\xi}}
\newcommand{\vnu}{\vect{\nu}}
\newcommand{\vlambda}{\vect{\lambda}}
\newcommand{\vzero}{\vect{0}}

% second-order tensors
\newcommand{\sig}{\tens{\sigma}}
\newcommand{\eps}{\tens{\varepsilon}}
\newcommand{\tI}{\tens{I}}               % identity
\newcommand{\tk}{\tens{k}}               % conductivity
\newcommand{\tA}{\tens{A}}               % coefficient field A(x)
\newcommand{\tM}{\tens{M}}               % polarization tensor (Ch. 3)
\newcommand{\tmu}{\tens{\mu}}            % moment tensor of the opening (Ch. 3)
\newcommand{\talpha}{\tens{\alpha}}      % thermal expansion

% fourth-order tensors
\newcommand{\bbC}{\tensf{C}}             % elastic moduli (Hooke)
\newcommand{\bbS}{\tensf{S}}             % compliances, S = C^{-1}
\newcommand{\bbI}{\tensf{I}}             % symmetric identity

% operators
\DeclareMathOperator{\dive}{div}
\DeclareMathOperator{\tr}{tr}
\DeclareMathOperator{\dev}{dev}
\DeclareMathOperator{\sym}{sym}
\DeclareMathOperator{\diag}{diag}
\DeclareMathOperator*{\argmin}{arg\,min}
\DeclareMathOperator*{\argmax}{arg\,max}
\newcommand{\T}{^{\!\top}}                % transpose
\newcommand{\dd}{\,\mathrm{d}}            % integration measure
\newcommand{\jump}[1]{[\![#1]\!]}
\newcommand{\norm}[1]{\lVert #1\rVert}
\newcommand{\scal}[2]{\langle #1,#2\rangle}
\newcommand{\dpr}{\!:\!}                   % double contraction (compact)

% boundary data (superscript D = "datum")
\newcommand{\uD}{u^{D}}
\newcommand{\pD}{\varphi^{D}}
\newcommand{\vuD}{\vu^{D}}
\newcommand{\vtD}{\vt^{D}}
\newcommand{\vfD}{\vf^{D}}

% Chapter 3 (reciprocity gap)
\newcommand{\RG}{\mathcal{R}}
\newcommand{\OG}{\Omega_\Gamma}

%======================================================================
%  Additions for the Cast3M notes
%
%  Everything above is shared verbatim with the "Nonlinear Problems in
%  Mechanics" notes, so that \vect, \tens, \sig, \eps, \bbC ... mean the
%  same thing in both books.  Below are the macros this book needs on top
%  of that, kept apart so the shared part can be updated by copying the
%  file across.
%======================================================================

% The old \tenq (4th-order, \pmb) is now an alias of the shared \tensf,
% so that fourth-order tensors are blackboard bold here too.
\newcommand{\tenq}[1]{\tensf{#1}}

\newcommand{\bfm}[1]{\mbox{\boldmath $#1$}}
\newcommand{\derp}[2]{\frac{\partial #1}{\partial #2}}

% grad and div written out, as in the Cast3M notices
\newcommand{\grad}{\text{grad}\,}
\newcommand{\ddiv}{\text{div}\,}
\newcommand{\mexp}{\text{exp}\,}

\newcommand{\bcddot}{\, \bfm{:} \,}
\newcommand{\bcdot}{\, \bfm{\cdot} \,}

\newcommand{\abs}[1]{\lvert #1\rvert}
\newcommand{\func}[1]{\mbox{\rm #1}}

% A gibiane operator or keyword in the running text.
\newcommand{\op}[1]{\texttt{#1}}

%  Marks a place where the original manuscript only referred to an equation
%  of the lecture notes.  Replace these as you fill them in; grep for them:
%      grep -n 'todoeq' ch_*.tex
\newcommand{\todoeq}[1]{%
  \ifmmode\text{\color{gray}\itshape[#1]}\else{\color{gray}\itshape[#1]}\fi}

%======================================================================
%  Index macros  (double-entry index, unchanged house style)
%
%  The index is a double-entry index: each gibiane operator is listed under
%  its own name followed by its english meaning, and again under the english
%  meaning followed by the operator.
%
%     \dindex{cerc[le]}{circle}  ->  cerc[le] (circle), 12
%                                    circle cerc[le], 12
%     \sindex{pasapas}           ->  pasapas, 34
%     \cindex{buckling}          ->  buckling, 43            (a concept)
%     \cindex[theta-method]{$\theta$-method}  -> sorts as "theta-method"
%
%  The sort keys ("#1 #2d@...") keep the operator entries next to each other.
%  NOTE on makeindex's three special characters, @ ! and |: in the VISIBLE
%  part of any of these macros they must be escaped with a double quote, as
%  in \sindex{"@excel1}.  An unescaped @ is read as a sort-key separator and
%  makeindex rejects the entry.
%======================================================================

\newcommand{\dindex}[2]{%
\index{#1 #2d@\texttt{#1} (#2)}%
\index{#2 #1@#2 \texttt{#1}}%
}

\newcommand{\sindex}[1]{%
\index{#1d@\texttt{#1}}%
}

\makeatletter
\newcommand{\cindex}[2][]{%
\def\cidx@sort{#1}%
\ifx\cidx@sort\@empty \index{#2@#2}\else \index{#1@#2}\fi
}
\makeatother

% Underscore inside a long Cast3M name, with a line break permitted after
% it.  Needed for names like Mooney\ul LRGTreloar\ul Cisaillementsimple,
% which are wider than the comment column of a worked example.
\newcommand{\ul}{\_\allowbreak}
